\tzpanchor{0806}
\tzpentryheader{0806}{LOGIT-SPHERICAL COORDINATE TRANSFORM}

\begingroup\small
\begingroup\scriptsize
\setlength{\tabcolsep}{4pt}
\renewcommand{\arraystretch}{0.92}
\noindent\begin{tabularx}{\linewidth}{@{}p{0.24\linewidth}p{0.36\linewidth}X@{}}
\textbf{UUID} & \multicolumn{2}{l@{}}{\tzpcode{e328a23d-a497-5d32-ba84-8a4d64ac6b4f}}\\
\textbf{PiID} & \tzpcode{9455346908302} & \textbf{TZPi}\quad \tzpcode{6}\quad \textbf{PiPE}\quad \tzpcode{813}\\
\textbf{RTEID} & \textbf{Poloidal $\theta$}\quad \tzpcode{3091.0448 rad} & \textbf{Toroidal $\phi$}\quad \tzpcode{05.0931 rad}\\
\textbf{TZPID} & \tzpcode{ID0806} & \textbf{Node Dependencies}\quad \tzpidref{0805}\\
\textbf{Registry} & \tzpcode{registry\_id\_0806} & \textbf{Scale}\quad transfer\\
\textbf{LOC Classification} & QA641 manifolds and topology & QC174.17 mathematical physics\\
\textbf{arXiv Classification} & Primary: math-ph & Cross-list: gr-qc, math.DG\\
\textbf{Primary Support} & Canonical Definition & Support Relation\\
\textbf{Closure Support} & Closure Relation & \\
\end{tabularx}\par
\endgroup
\endgroup

\tzpsection{Canonical Statement}
The entry records logit-spherical coordinate transform through a normalized Nonary TRAWIN operator chain.

\tzpsection{Canonical Equation}
\[
\mathcal{O}_{\mathrm{id0806}}:\mathcal{X}_{\mathrm{id0806}}\longrightarrow\mathcal{Y}_{\mathrm{id0806}}
\]
\[
\mathcal{N}_{\mathrm{TRAWIN}}(\mathcal{O}_{\mathrm{id0806}})=\mathcal{O}_{\mathrm{id0806}}
\]
\[
\mathcal{O}_{\mathrm{id0806}}(x)\in\mathcal{Y}_{\mathrm{adm}}
\]

\begin{minipage}[t]{0.49\linewidth}
\tzpsection{Canonical Symbol Key}
\begingroup
\small
\raggedright
\textbf{Source equation symbols.}\par
\begin{itemize}[leftmargin=1.35em,itemsep=1pt,topsep=1pt,parsep=0pt]
    \item $\mathcal{O}_{\mathrm{id0806}}$: canonical operator or carrier map for the entry.
    \item $\mathcal{X}_{\mathrm{id0806}}$: source domain or input state space for the entry map.
    \item $\mathcal{Y}_{\mathrm{id0806}}$: target domain or output state space for the entry map.
    \item $\mathcal{N}_{\mathrm{TRAWIN}}$: TRAWIN normalization operator applied to the entry relation.
    \item $\mathcal{Y}_{\mathrm{adm}}$: admissible output space selected by the source equation.
\end{itemize}
\endgroup
\end{minipage}\hfill
\begin{minipage}[t]{0.47\linewidth}
\tzpsection{Wolfram Form}
\begingroup
\scriptsize
\raggedright
\textbf{Source Wolfram model.}\par
\begin{Verbatim}[fontsize=\tiny,breaklines=true,breakanywhere=true]
(* TZPID ID0806 generated source Wolfram model *)
ClearAll[K0806, Cr0806, T, R, A, W, I, N];
eq0806$1 = HoldForm[Oid0806[Xid0806] -> Yid0806];
eq0806$2 = HoldForm[NTRAWIN[OId0806] == OId0806];
eq0806$3 = HoldForm[Element[OId0806[x], YAdm]];

K0806 = HoldForm[{eq0806$1, eq0806$2, eq0806$3}];

N = "normalized TRAWIN closure object for LOGIT-SPHERICAL COORDINATE TRANSFORM";
I = "LOGIT-SPHERICAL COORDINATE TRANSFORM integration/comparison layer";
W = "LOGIT-SPHERICAL COORDINATE TRANSFORM wave or operator carrier";
A = "LOGIT-SPHERICAL COORDINATE TRANSFORM admissible action/amplitude condition";
R = "LOGIT-SPHERICAL COORDINATE TRANSFORM rotation/curl preservation layer";
T = "LOGIT-SPHERICAL COORDINATE TRANSFORM local source condition";

Cr0806[expr_] := HoldForm[N[I[W[A[R[T[expr]]]]]]];

Result0806 = Cr0806[K0806];
\end{Verbatim}
\endgroup
\end{minipage}

\par\vspace{0.45\baselineskip}
\noindent\begin{minipage}[t]{\linewidth}
\begingroup
\small
\raggedright
\textbf{TRAWIN Operator Key.}\par
\noindent\textbf{Time/localization operator:} $\mathsf{T}\!\left[\mathrm{local}\right]$ LOGIT-SPHERICAL COORDINATE TRANSFORM local source condition.\par
\noindent\textbf{Rotation/curl operator:} $\mathsf{R}\!\left[\nabla\times\right]$ LOGIT-SPHERICAL COORDINATE TRANSFORM rotation/curl preservation layer.\par
\noindent\textbf{Action/amplitude operator:} $\mathsf{A}\!\left[\mathrm{admissible}\right]$ LOGIT-SPHERICAL COORDINATE TRANSFORM admissible action/amplitude condition.\par
\noindent\textbf{Wave/d'Alembertian operator:} $\mathsf{W}\!\left[\Box\right]$ LOGIT-SPHERICAL COORDINATE TRANSFORM wave or operator carrier.\par
\noindent\textbf{Integration operator:} $\mathsf{I}\!\left[\int\right]$ LOGIT-SPHERICAL COORDINATE TRANSFORM integration/comparison layer.\par
\noindent\textbf{Normalization operator:} $\mathsf{N}\!\left[\mathrm{normalized}\right]$ normalized TRAWIN closure object for LOGIT-SPHERICAL COORDINATE TRANSFORM.\par
\endgroup
\end{minipage}

\clearpage
\noindent\textbf{[ID 0806] LOGIT-SPHERICAL COORDINATE TRANSFORM}\par
\vspace{0.35em}
\tzpsection{Derivation Layer}
\paragraph{Primary Equation.}
\[
\mathcal{O}_{\mathrm{id0806}}:\mathcal{X}_{\mathrm{id0806}}\longrightarrow \mathcal{Y}_{\mathrm{id0806}},\qquad
\mathcal{N}_{\mathrm{TRAWIN}}(\mathcal{O}_{\mathrm{id0806}})=\mathcal{O}_{\mathrm{id0806}},\qquad
\mathcal{O}_{\mathrm{id0806}}(x)\in\mathcal{Y}_{\mathrm{adm}}.
\]

\paragraph{Primary Derivation.} Let $x\in \mathcal{X}_{\mathrm{id0806}}$ be the local coordinate or field sample.  The entry transform is\[
y=\mathcal{O}_{\mathrm{id0806}}(x),\qquad \mathcal{O}_{\mathrm{id0806}}:\mathcal{X}_{\mathrm{id0806}}\longrightarrow \mathcal{Y}_{\mathrm{id0806}}.
\]

Its differential transport is represented by the id-local Jacobian\[
J_{0806}(x)=D\mathcal{O}_{\mathrm{id0806}}\big|_x,\qquad \det J_{0806}(x)\neq 0
\]

Compatibility with the Nonary chain is enforced by\[
\mathcal{N}_{\mathrm{TRAWIN}}(\mathcal{O}_{\mathrm{id0806}})=\mathcal{O}_{\mathrm{id0806}},\qquad y\in\mathcal{Y}_{\mathrm{adm}}.
\]

Thus the Nonary derivation for [ID 0806] is the three-stage passage
\[
\mathcal{X}_{\mathrm{id0806}}\xrightarrow{\,\mathcal{O}_{\mathrm{id0806}}\,} \mathcal{Y}_{\mathrm{id0806}}\xrightarrow{\,\mathcal{N}_{\mathrm{TRAWIN}}\,} \mathcal{Y}_{\mathrm{adm}},
\]
with the entry title fixing the meaning of the carrier and the TRAWIN normalization fixing the admissibility test.

\paragraph{Interpretation.} LOGIT-SPHERICAL COORDINATE TRANSFORM is read as a transport rule: it converts local data into an admissible Nonary coordinate while preserving enough differential structure for the following entry to use it.

\par\vspace{0.35em}
\noindent\textbf{Derivable Consequences.}\quad
\begingroup
\footnotesize
\raggedright
The canonical equation anchors LOGIT-SPHERICAL COORDINATE TRANSFORM by connecting the source condition to the derivation layer and proof certificate.
\par\endgroup

\tzpsection{Equation Support}
\noindent\textbf{Canonical Support Relation}\par
\[
\mathcal{O}_{\mathrm{id0806}}:\mathcal{X}_{\mathrm{id0806}}\longrightarrow\mathcal{Y}_{\mathrm{id0806}}
\]
\[
\mathcal{N}_{\mathrm{TRAWIN}}(\mathcal{O}_{\mathrm{id0806}})=\mathcal{O}_{\mathrm{id0806}}
\]
\[
\mathcal{O}_{\mathrm{id0806}}(x)\in\mathcal{Y}_{\mathrm{adm}}
\]
\noindent The support equation repeats the canonical relation so the derivation page remains self-contained.\par

\clearpage
\noindent\textbf{[ID 0806] LOGIT-SPHERICAL COORDINATE TRANSFORM}\par
\vspace{0.35em}
\tzpsection{Dictionary}
\begingroup\small
\tzpsection{Definition}
LOGIT-SPHERICAL COORDINATE TRANSFORM is a registry definition in Resonance and Phase Locking with secondary pressure from Biological and Biomolecular Analogies.

% BEGIN TZP CONTEXTUAL MEANING
\tzpsection{Contextual Meaning}
\begingroup\footnotesize\setlength{\emergencystretch}{3em}\sloppy
\textbf{Formal statement:} mathcal E encoder (M) = text CantorPair [text Logit [text Mod [hat n * text diag (pi digits ), 1]]]. \textbf{Contextual meaning:} The entry reads LOGIT-SPHERICAL COORDINATE TRANSFORM as a controlled technical headword whose equation layer, proof route, and registry role are interpreted together. \textbf{Derivable consequences:} The entry now reads through actual sourced equations, so its local arc is carried by explicit mathematical relations rather than a uniform quaternary certificate record shell. \textbf{Lemma support:} Canonical Definition; Support Relation; Closure Relation.\par
\endgroup
% END TZP CONTEXTUAL MEANING

\tzpsection{Technical Interpretation}
Technically, LOGIT-SPHERICAL COORDINATE TRANSFORM is read through the local equation chain and the TRAWIN normalization recorded on the entry page.

\tzpsection{Equation Grounding}
Equation anchors: mathcal E encoder (M) = CantorPair [ Logit [ Mod [ hat n dot diag (pi digits ), 1]]]; mathcal J sub out = T sub TZP ( K sub TZP ); and mathcal A sub adm ( J sub out )=1. Proof route: show that the stated equations form a coherent progression for the entry and remain compatible under the TRAWIN ordering. Formalization route: prove that the final closure relation follows from the preceding entry-specific equations.

\tzpsection{Interpretive Role}
The entry now reads through actual sourced equations, so its local arc is carried by explicit mathematical relations rather than a uniform quaternary certificate record shell.
\endgroup

\tzpsection{TZP Thesaurus}
\begin{enumerate}[leftmargin=1.6em,itemsep=6pt,topsep=4pt]
\item \textbf{Infinite-Plane to Globe Wrapping}: The dizzying mathematical trick of taking a completely flat, never-ending sheet of data and perfectly, seamlessly wrapping it around the surface of a 3D sphere.
\item \textbf{Probability-Space Squashing}: The compression algorithm that forces impossible, infinite odds to fit neatly between the physical boundaries of 0 and 100\%.
\item \textbf{Non-Linear Polar Mapping}: The mapping system that wildly stretches the grid lines near the poles to ensure the math doesn't break down when the energy flows over the top of the reactor.
\end{enumerate}

\tzpsection{Encyclopedia Mathematica}
\begingroup\footnotesize\sloppy
The visual plate below is reserved for the source-specific equation and progression graphic for [ID 0806]. It is included only as a companion to the canonical equation, not as a substitute for the formal text.
\par\endgroup
\begingroup\footnotesize\itshape
Visual plate pending source-specific approval for [ID 0806].
\par\endgroup

\clearpage
\begingroup\tzpsupportsetup
\noindent\textbf{\fontsize{10.8pt}{12.8pt}\selectfont [ID 0806] Constructive Logic and Proof Certificate}\par
\noindent\textbf{\fontsize{10pt}{12pt}\selectfont LOGIT-SPHERICAL COORDINATE TRANSFORM}\par
\vspace{0.45em}
\begingroup
\footnotesize
\setlength{\parskip}{0.22em}
\setlength{\parindent}{0pt}
\noindent\textbf{Curry--Howard--Lambek Interpretation}\par
Under the Curry--Howard--Lambek correspondence, this registry entry is read as a typed proof
object: the stated source condition supplies the proposition, the derivation supplies the
morphism, and the proof certificate records the machine handles needed to check the obligation
in the formal registry.

\vspace{0.45em}
\noindent\textbf{CHL Proof}\par
\textbf{Statement:} LOGIT-SPHERICAL COORDINATE TRANSFORM is read as a source-backed proof obligation.

\textbf{Logic Validation:}\\
\textbf{Proposition:} ``LOGIT-SPHERICAL COORDINATE TRANSFORM is admissible under the named source equation and derivation.''\\
\textbf{Proof:} The canonical statement supplies the proposition, the derivation supplies the witness, and the TRAWIN operator chain records the normalization path.

\textbf{VERDICT:} \ensuremath{\checkmark}\ \textbf{TRUE} --- conditional registry validation

\textbf{Formal Status:} \ensuremath{\checkmark}\ THEORETICAL -- source-backed CHL proof obligation

\textbf{Physical Status:} THEORETICAL -- requires model-specific empirical interpretation
\par\vspace{0.45em}
\noindent{\tiny\textbf{Isabelle/HOL.} \tzpplaincode{registry\_number\_matches, equation\_count\_matches}}\par
\vfill
\begingroup\footnotesize\itshape
Proof visual pending source-specific approval for [ID 0806].
\par\endgroup
\vfill
\endgroup
\endgroup
